% Erweiterungen zur Beschreibung des Aufbaus der Arbeit

\noindent\textbf{Kapitel 7}\quad
Im siebten Kapitel werden die Messergebnisse der drei untersuchten Zephyr-Konfigurationen Minimal, Logging und IoT-Stack dargestellt und miteinander verglichen. Zuerst werden die Ergebnisse der beiden Kontextwechselvarianten, der Semaphor-Weckung eines h{\"o}herprioren Threads und der kooperativen {\"U}bergabe durch \lstinline[style=codeinline]|k_yield()|, betrachtet. Anschlie{\ss}end folgen die Messwerte der Interrupt-Latenz und des Timer-Jitters. Danach werden die Ergebnisse zur Geschwindigkeit der Heap-Allokation und -Freigabe und zum Verhalten des Speichers nach einem gezielt erzeugten Fragmentierungsmuster dargestellt. Abschlie{\ss}end wird der statische Flash- und RAM-Bedarf der drei Konfigurationsprofile miteinander verglichen. Das Kapitel konzentriert sich dabei auf die deskriptive Darstellung der ermittelten Werte. Ihre Ursachen und Bedeutung werden daher erst im darauffolgenden Kapitel diskutiert.

\noindent\textbf{Kapitel 8}\quad
Im achten Kapitel werden die davor dargestellten Messergebnisse interpretiert und dann im Hinblick auf die Forschungsfrage eingeordnet. Zuerst wird untersucht, welchen Einfluss die Aktivierung des Logging-Subsystems auf die gemessenen Laufzeiten und den Speicherbedarf hat. Anschlie{\ss}end wird der Einfluss des IoT-Stacks mit OpenThread, CoAP und DTLS/PSK betrachtet. Dabei wird zwischen dem Ressourcenbedarf der eingebundenen Subsysteme und den Kosten ihrer Nutzung unterschieden. Das liegt daran, dass w{\"a}hrend der Mikrobenchmarks keine Log-Nachrichten, kein aktiver Netzwerkverkehr und kein DTLS-Handshake erzeugt werden. Danach werden die erkennbaren Zielkonflikte zwischen Funktionsumfang, Laufzeitverhalten und Speicherverbrauch herausgearbeitet. Zuletzt werden die eigenen Ergebnisse mit den in Kapitel 3 vorgestellten Untersuchungen und Referenzwerten aus der Literatur verglichen.

\noindent\textbf{Kapitel 9}\quad
Im neunten Kapitel wird die Validit{\"a}t der durchgef{\"u}hrten Untersuchung bewertet. Zuerst wird die interne Validit{\"a}t betrachtet, also wie sehr die beobachteten Unterschiede auf die verschiedenen Zephyr-Konfigurationen zur{\"u}ckgef{\"u}hrt werden k{\"o}nnen. Anschlie{\ss}end wird bei der externen Validit{\"a}t untersucht, wie weit die Ergebnisse auf andere Hardwareplattformen, Zephyr-Versionen und Konfigurationen {\"u}bertragbar sind. Dar{\"u}ber hinaus wird der Messaufwand, der durch das Auslesen des DWT-Zyklenz{\"a}hlers entsteht, inklusive seiner Kalibrierung behandelt. Zuletzt werden die Limitierungen des Versuchsaufbaus dargestellt. Dazu geh{\"o}ren die Durchf{\"u}hrung jeweils einer zusammenh{\"a}ngenden Messreihe pro Messzelle, m{\"o}gliche Einfl{\"u}sse des Flash-Layouts, die begrenzte Aufl{\"o}sung des RTC-Systemtimers und die fehlende aktive Logging- und Netzwerkbelastung w{\"a}hrend der Benchmarks.

\noindent\textbf{Kapitel 10}\quad
Im zehnten Kapitel werden die relevantesten Erkenntnisse der Untersuchung zusammengefasst und die in Kapitel 1 formulierte Forschungsfrage inklusive ihrer Unterfragen beantwortet. Dabei werden die Auswirkungen der drei Konfigurationsprofile auf die Kontextwechsel- und Interrupt-Latenz, den Timer-Jitter, die Heap-Operationen, die Heap-Fragmentierung und den statischen Flash- und RAM-Bedarf zusammengefasst. Im Mittelpunkt steht dabei die Erkenntnis, dass weitere Zephyr-Subsysteme vor allem den Speicherbedarf deutlich erh{\"o}hen. Auf der anderen Seite ist ihr Einfluss auf die untersuchten kurzen Laufzeitpfade geringer und folgt nicht durchgehend einer monotonen Entwicklung. Au{\ss}erdem werden die praktischen Erkenntnisse aus der Inbetriebnahme des IoT-Profils zusammengefasst.

\noindent\textbf{Kapitel 11}\quad
Im elften Kapitel werden m{\"o}gliche Erweiterungen und weiterf{\"u}hrende Untersuchungen betrachtet. Dazu geh{\"o}ren mehrere voneinander unabh{\"a}ngige Build-, Flash- und Neustartvorg{\"a}nge pro Messzelle und die Erfassung der kompletten Einzelmesswerte. Dar{\"u}ber hinaus k{\"o}nnen zuk{\"u}nftige Messungen aktive Logging-Lasten, einen OpenThread-Netzwerkbeitritt, CoAP-Kommunikation und DTLS-Handshakes einbeziehen. Weitere Erweiterungen bestehen in einer h{\"o}her aufl{\"o}senden Untersuchung des Timer-Jitters, einer gezielten Analyse von m{\"o}glichen Flash-Layout-Effekten und einer Messung der Stackauslastung. Zuletzt wird die {\"U}bertragung des Versuchsaufbaus auf weitere Mikrocontrollerarchitekturen, Zephyr-Versionen, Konfigurationsprofile und andere Echtzeitbetriebssysteme als Bausteine f{\"u}r zuk{\"u}nftige Arbeiten vorgeschlagen.
